\uuid{0QYo}
\titre{Solution générale/raccord}
\auteur{hanani ; mzari}
\organisation{exo7}
\licence{CC BY-NC-SA 4.0}
\datecreate{2020-01-20}
\niveau{L1}
\module{Analyse}
\chapitre{Equation différentielle}
\sousChapitre{Résolution d'équation différentielle du premier ordre}
\difficulte{3}
\theme{Equations du premier ordre, solution générale, raccord}
\exo7qcmid{lille-2.equations-differentielles.017}

\contenu{
\begin{question}
\qclassification{225.01}
\qtags{motcle=solution générale/raccord}
On considère l'équation différentielle $(E)$ : $xy'-2y=x$. Parmi les affirmations suivantes, cocher celles qui sont vraies :
\begin{answers}
\good{La solution générale de $(E)$ sur $]0,+\infty[$ est : $y=kx^2-x$, $k\in \Rr$.}
\good{La solution générale de $(E)$ sur $]-\infty,0[$ est : $y=kx^2-x$, $k\in \Rr$.}
\bad{Les solutions de $(E)$ sur $\Rr$ sont les fonctions $y(x)=kx^2-x$, où $k\in \Rr$.}
\bad{$(E)$ possède une seule solution sur $\Rr$.}
\end{answers}
\begin{explanations}
La solution générale de l'équation homogène est : $y=kx^2$, $k\in \Rr$, et $y_0=-x$ est une solution particulière. Si $y$ est une solution de $(E)$ sur $\Rr$, alors $y$ est une solution de $(E)$ sur $]0,+\infty[$ et $y$ est une solution de $(E)$ sur $]-\infty,0[$. De plus $y$ est continue en $0$. D'où
$$y(x)=\left\{\begin{array}{lll}k_1x^2-x&\mbox{si}&x\in ]0,+\infty [ \\k_2x^2-x&\mbox{si}&x\in ]-\infty ,0[\\ 0&\mbox{si}&x=0. \end{array}\right.$$
Avec $k_1,k_2\in \Rr$. Une telle fonction est dérivable en $0$. réciproquement, toute fonction de la forme ci-dessus est une solution de $(E)$ sur $\Rr$.
\end{explanations}
\end{question}
}
